\midsectiontitle{O Alento do Aventureiro}

\section*{Tréguas em Meio ao Caos}

No domínio implacável do Pináculo, o aço cansa e o espírito enfraquece. O descanso não é apenas um luxo, mas uma necessidade vital para aqueles que pretendem ver o sol nascer mais uma vez. Em \textit{Overgrown}, o repouso é dividido entre a rápida recuperação de fôlego e o profundo restabelecimento das forças.

\section*{Tipos de Descanso}

\subsection*{Pausa Breve (1 a 2 horas)}
Uma parada estratégica para estancar sangramentos e recuperar o fôlego. Durante uma Pausa Breve, o aventureiro deve rolar \textbf{2D8} e decidir qual resultado atribuir aos seus Pontos de Vida e qual aos seus Pontos de Energia Arcana. 

\begin{rpgboxtips}{Vulnerabilidade}
Enquanto descansam, os aventureiros tornam-se obrigatoriamente \link{status:desprevenido}{Desprevenidos}. Para evitar emboscadas fatais, um membro do grupo pode assumir a \link{status:vigilia}{Vigília}. Aquele que fica de guarda não recupera atributos, mas garante que o grupo não seja pego de surpresa por predadores ou patrulhas inimigas.
\end{rpgboxtips}

\subsection*{Descanso Pleno (Noite Completa)}
Um sono profundo, acompanhado de alimentação e hidratação adequadas. Este descanso exige um ambiente minimamente seguro e o consumo de rações ou refeições preparadas. A recuperação é massiva:
\begin{itemize}
    \item \textbf{Vitalidade:} Recupera 1D20 + FOR Pontos de Vida.
    \item \textbf{Energia Arcana:} Recupera 1D20 + WIS Pontos de Energia.
\end{itemize}

\section*{A Arte do Banquete}

Embora as rações de viagem mantenham um aventureiro de pé, apenas uma refeição preparada com maestria pode elevar o potencial de um aventureiro. Através da perícia \link{per:chef}{Chef*}, é possível transformar ingredientes crus em iguarias que potencializam os descansos.

Abaixo estão os níveis de complexidade para o preparo ou compra de refeições e seus respectivos bônus (cumulativos com a rolagem natural do descanso):

\begin{center}
    \small 
    \renewcommand{\arraystretch}{1.4}
    \begin{tabularx}{\columnwidth}{@{} l c r r @{}}
        \toprule
        \rowcolor{gray!15} 
        \textbf{Qualidade} & \textbf{Dificuldade} & \textbf{Breve} & \textbf{Pleno} \\ 
        \midrule
        Básica      & 10 & +1D4  & +1D6  \\
        Simples     & 15 & +1D6  & +1D8  \\
        Mediana     & 20 & +1D8  & +1D10 \\
        Avançada    & 25 & +1D10 & +1D12 \\
        Mestre      & 30 & +1D12 & +1D20 \\
        \rowcolor{overgrownRed!10}
        Suprema     & 35 & +1D20 & +2D20 \\
        \bottomrule
    \end{tabularx}
\end{center}

\subsection*{Níveis de Gastronomia}

\begin{itemize}
    \item \textbf{Básica:} Um cozido simples, feito com o que há disponível. Exige pouco esforço, mas aquece o peito.
    \item \textbf{Simples:} Uso de temperos comuns e técnica de cocção correta. Proporciona um conforto notável.
    \item \textbf{Mediana:} Requer um conhecimento decente de ervas e cortes, extraindo o melhor das propriedades dos alimentos.
    \item \textbf{Avançada:} Utiliza especiarias de múltiplos locais. O sabor é tão intenso que revigora as forças instantaneamente.
    \item \textbf{Mestre:} Um equilíbrio perfeito entre temperatura e natureza dos ingredientes. O prato parece brilhar com energia.
    \item \textbf{Suprema:} Uma obra-prima culinária. Dizem que tais pratos podem fazer um homem à beira da morte se sentir um deus por alguns instantes.
\end{itemize}

\begin{rpgboxtips}{Bebidas e Tabernas}
Refeições especiais também podem ser compradas em tavernas. O mestre define o custo baseado na raridade dos ingredientes e no nível de pericia de \link{per:chef}{Chef} do cozinheiro local.
\end{rpgboxtips}